\section*{Bab \ref{ch:b2:quadratic} --- Bentuk Kuadratik}

\begin{solution}{exo:b2:quadratic:1}
$q_1 = (x + 2y)^2 - 3y^2$: ranknya $2$ dan signaturnya $(1, 1)$ (yakni
bentuk bertipe hiperbola).

$q_2$: lengkapkan kuadratnya dalam $x$: $q_2 = (x + y)^2 + y^2 + 2yz +
3z^2 = (x+y)^2 + (y + z)^2 + 2z^2$: jadi ranknya $3$ dan signaturnya $(3, 0)$
--- sehingga definit positif.
\end{solution}

\begin{solution}{exo:b2:quadratic:2}
\omterm{def:b2:reduction:eigen}{Nilai eigennya} $2$ (pada $\operatorname{Vect}(1,1,1)$) dan $-1$ (pada
bidang $x + y + z = 0$). Ortonormalkan: $e_1 =
\frac{1}{\sqrt3}(1,1,1)$; lalu pada bidangnya, Gram--Schmidt atas $(1,-1,0),
(1,0,-1)$ memberi $e_2 = \frac{1}{\sqrt2}(1,-1,0)$ dan $e_3 =
\frac{1}{\sqrt6}(1,1,-2)$. Maka $P = (e_1\ e_2\ e_3)$ bersifat ortogonal
dengan $P^{\mathsf T}AP = \operatorname{diag}(2,-1,-1)$.

Bentuk $q = 2xy + 2yz + 2zx$ bermatriks $A$: jadi dalam koordinat
terputarnya $q = 2X^2 - Y^2 - Z^2$ --- yakni sumbu utamanya; signaturnya
$(1,2)$, sepakat dengan \cref{ex:b2:quadratic:gaussexample}
(sebab bentuknya sama!).
\end{solution}

\begin{solution}{exo:b2:quadratic:3}
$u^{**} = u$: sebab $\langle u^{**}x, y\rangle = \langle x, u^*y\rangle =
\langle ux, y\rangle$ untuk setiap $y$. Lalu $(uv)^* = v^*u^*$: sebab $\langle
uvx, y\rangle = \langle vx, u^*y\rangle = \langle x,
v^*u^*y\rangle$. Kernelnya: $y \in \ker u^* \iff \langle x, u^*y
\rangle = 0\ \forall x \iff \langle u(x), y\rangle = 0\ \forall x
\iff y \perp \operatorname{im} u$. Ranknya: $\dim\ker u^* = n -
\operatorname{rk} u$ (lewat komplemen ortogonalnya), jadi
$\operatorname{rk} u^* = \operatorname{rk} u$ menurut rank--nulitas ---
yakni awatara Euklides bagi teorema rank \omterm{def:b2:linalg:transpose}{transpos}.
\end{solution}

\begin{solution}{exo:b2:quadratic:4}
Secara spektral: $A = PDP^{\mathsf T}$ dengan $D$ diagonal berentri
$\lambda_i$ yang memenuhi $\lambda_i^3 = \lambda_i$: jadi $\lambda_i \in
\{-1, 0, 1\}$. Maka $A^2 = PD^2P^{\mathsf T}$ dengan $D^2$ diagonal
berentri $0/1$: jadi $A^2$ \omterm{def:b2:quadratic:adjoint}{simetrik} sekaligus idempoten
(sebab $(A^2)^2 = A^4 = A\cdot A^3 = A^2$) --- dan \omterm{def:b2:quadratic:adjoint}{simetrik} idempoten
$=$ proyeksi ortogonal (ia proyeksi ke
$\ker(A^2 - I) = \ker(A-I)\oplus\ker(A+I)$ sepanjang $\ker A$, dan
keduanya ortogonal menurut teorema spektral).

Secara umum, $P(A) = P\!\left(\text{diag}\right)$: jadi $P(A)$ punya
\omterm{def:b2:reduction:eigen}{vektor eigen} yang sama dengan \omterm{def:b2:reduction:eigen}{nilai eigen} $P(\lambda_i)$ --- yakni
``pemetaan spektral'' pada aras yang \omterm{def:b2:reduction:diag}{dapat didiagonalkan}.
\end{solution}

\begin{solution}{exo:b2:quadratic:5}
Tertutup: $O(n) = g^{-1}(\{I\})$ untuk $g(P) =
P^{\mathsf T}P$ yang \omterm{def:b2:metric:continuity}{kontinu} (sebab entrinya polinomial). Terbatas: sebab tiap
kolom $P \in O(n)$ vektor satuan, jadi semua entrinya terletak di
$\intcc{-1}{1}$. Tertutup dan terbatas dalam $\mathcal{M}_n(\R) \simeq
\R^{n^2}$: sehingga \omterm{def:b2:metric:compact}{kompak}
(\cref{thm:b2:metric:compactprops}~(2)).

Tak \omterm{def:b2:metric:connected}{terhubung}: sebab $\det$ mengambil dua nilai $\pm1$ pada $O(n)$, dan
sebuah surjeksi \omterm{def:b2:metric:continuity}{kontinu} ke $\{-1, 1\}$ memecah ruangnya
(lewat hujah \cref{ex:b2:metric:glnr}).
\end{solution}

\begin{solution}{exo:b2:quadratic:6}
\emph{Keberadaannya:} $A = PDP^{\mathsf T}$ dengan $D =
\operatorname{diag}(\lambda_i)$ dan $\lambda_i \geq 0$; tetapkan $B =
P\sqrt D P^{\mathsf T}$ dengan $\sqrt D =
\operatorname{diag}(\sqrt{\lambda_i})$: ia \omterm{def:b2:quadratic:adjoint}{simetrik}, semidefinit
positif, dan $B^2 = A$.

\emph{Ketunggalannya:} misalkan $B$ \omterm{def:b2:quadratic:adjoint}{simetrik} semidefinit positif dengan
$B^2 = A$. Maka $B$ komut dengan $A$; sehingga $B$ memelihara tiap ruang eigen
$E_\lambda(A)$ (sebab untuk $Ax = \lambda x$: $A(Bx) = BAx = \lambda Bx$).
Pada $E_\lambda(A)$, pembatasan $B$ bersifat \omterm{def:b2:quadratic:adjoint}{simetrik} semidefinit positif
dengan kuadrat $\lambda\,\mathrm{id}$; jadi \omterm{def:b2:reduction:eigen}{nilai eigennya} $\mu$ memenuhi $\mu^2
= \lambda$ dan $\mu \geq 0$: sehingga $\mu = \sqrt\lambda$ --- jadi
pembatasannya, yang \omterm{def:b2:reduction:diag}{dapat didiagonalkan} dengan satu-satunya \omterm{def:b2:reduction:eigen}{nilai eigen}
$\sqrt\lambda$, \emph{adalah} $\sqrt\lambda\,\mathrm{id}$. Karena $E =
\bigoplus E_\lambda(A)$, maka $B$ tertentukan: $B = \sqrt A$.
\end{solution}

\begin{solution}{exo:b2:quadratic:7}
Dalam basis eigen ortonormal, $\norm{Ax}_2^2 = \sum \lambda_i^2
x_i^2 \leq (\max_i \lambda_i^2)\norm x_2^2$, dengan kesamaan pada \omterm{def:b2:reduction:eigen}{vektor
eigen} yang bersesuaian: jadi $\vertiii A_2 = \max\abs{\lambda_i}$.
Untuk matriks yang diberikan: \omterm{def:b2:reduction:eigen}{nilai eigennya} $3, -1$
(yakni kembaran \cref{ex:b2:quadratic:spectralexample}), jadi
$\vertiii A_2 = 3$.
\end{solution}

\begin{solution}{exo:b2:quadratic:8}
($\Rightarrow$) Blok kiri atas berukuran $k \times k$, sebut $A_k$, adalah
matriks pembatasan bentuk (yang definit) itu ke rentang $k$ vektor
basis pertamanya: ia definit positif, jadi \omterm{def:b2:reduction:eigen}{nilai eigennya}
positif dan $\Delta_k = \det A_k > 0$.

($\Leftarrow$) Secara induktif atas $n$; kasus $n = 1$ jelas. Andaikan semua
$\Delta_k > 0$. Menurut induksinya, $A_{n-1}$ definit positif: jadi
bentuk $q$ yang dibatasi ke $F = \operatorname{Vect}(e_1, \dots,
e_{n-1})$ bersifat definit. Diagonalkan $q|_F$ (lewat Gauss): dengan koordinat
$y_1, \dots, y_{n-1}$ berlaku $q|_F = \sum y_i^2$. Lalu di ruang penuhnya,
setelah melengkapkan kuadratnya dalam peubah terakhirnya,
\[
q = \sum_{i=1}^{n-1} \bigl(y_i + c_i x_n\bigr)^2 + c\,x_n^2
\]
untuk konstanta yang sesuai (kumpulkan suku silangnya ke dalam kuadratnya).
Reduksinya memamerkan signatur $(n-1 + \epsilon, \cdot)$ dengan
$\epsilon$ menyatakan sumbangan tanda $c$; dan \omterm{def:b2:linalg:det}{determinannya} memelihara
tanda hasil kali koefisien diagonalnya di bawah
\omterm{def:b2:quadratic:def}{kekongruenan} (sebab $\det(P^{\mathsf T}AP) = (\det P)^2\det A$): jadi $\Delta_n
> 0$ memaksa $c > 0$. Karena itu $q$ adalah jumlah $n$ kuadrat
bentuk yang bebas: sehingga definit positif.
\end{solution}

\begin{solution}{exo:b2:quadratic:9}
Misalkan $(e_1, \dots, e_n)$ basis eigen ortonormal bagi $\lambda_1
\geq \dots \geq \lambda_n$.

\emph{$\lambda_2 \leq$ min-maksnya:} untuk sembarang hiperbidang $H$,
ruang berdimensi $2$, yakni $V = \operatorname{Vect}(e_1, e_2)$, memenuhi
$\dim(H \cap V) \geq 1$ (menurut Grassmann): jadi pilihlah $x \in H \cap V$
yang satuan, $x = ae_1 + be_2$ dengan $a^2 + b^2 = 1$:
\[
\langle u(x), x\rangle = \lambda_1 a^2 + \lambda_2 b^2 \geq
\lambda_2 :
\]
jadi maksimum tiap hiperbidangnya $\geq \lambda_2$.

\emph{Untuk $\geq$:} bagi $H = e_1^{\perp}$, setiap $x = \sum_{i\geq2}
x_ie_i \in H$ yang satuan punya $\langle u(x), x\rangle = \sum_{i \geq 2}
\lambda_i x_i^2 \leq \lambda_2$, yang tercapai di $e_2$: jadi
maksimum hiperbidang ini persis $\lambda_2$. Karena itu minimum atas $H$
adalah $\lambda_2$.
\end{solution}

\begin{solution}{exo:b2:quadratic:10}
\emph{Jalan aljabarnya:} $2q = \bigl(\sum x_i\bigr)^2 - \sum
x_i^2$. Pada hiperbidang $H : \sum x_i = 0$ (yang berdimensi $n -
1$), $q = -\frac12\sum x_i^2$ bersifat definit negatif; sedangkan pada garis
$\R(1, \dots, 1)$ berlaku $q(t, \dots, t) = \binom n2 t^2 > 0$. Sebuah
subruang tempat $q$ definit positif memotong $H$ secara trivial, jadi
berdimensi $\leq 1$: sehingga menurut Sylvester
(\cref{thm:b2:quadratic:sylvester}), $s = 1$, dan $t \geq n-1$
berkat $H$; lalu rank $\leq n$ memaksa signaturnya $(1, n-1)$ dan ranknya $n$.

\emph{Jalan spektralnya:} matriksnya $\frac12(J - I)$; sedangkan $J$ punya
\omterm{def:b2:reduction:eigen}{nilai eigen} $n$ (pada $(1,\dots,1)$) dan $0$ (pada $H$), jadi
$\frac12(J-I)$ bernilai eigen $\frac{n-1}{2}$ (sekali) dan
$-\frac12$ (sebanyak $n-1$ kali): jadi satu positif dan $n-1$ negatif ---
signatur yang sama, menurut \cref{cor:b2:quadratic:principalaxes}.
\end{solution}

\begin{solution}{exo:b2:quadratic:11}
Tulislah $B = C^{\mathsf T}C$ (\cref{exo:b2:quadratic:6} lewat $C =
\sqrt B$). Maka, dengan $c_1, \dots, c_n$ menyatakan kolom
$C^{\mathsf T}$:
\[
\operatorname{tr}(AB) = \operatorname{tr}(AC^{\mathsf T}C)
= \operatorname{tr}(CAC^{\mathsf T})
= \sum_{i=1}^n \langle A c_i, c_i\rangle .
\]
Menurut \cref{cor:b2:quadratic:principalaxes}~(2), tiap sukunya terletak
antara $\lambda_{\min}(A)\norm{c_i}^2$ dan
$\lambda_{\max}(A)\norm{c_i}^2$, dan $\sum\norm{c_i}^2 =
\operatorname{tr}(C^{\mathsf T}C)^{\vphantom1} =
\operatorname{tr} B$: jadi ketaksamaan gandanya menyusul. Bila $A$
juga semidefinit positif, maka $\lambda_{\min}(A) \geq 0$ sehingga
$\operatorname{tr}(AB) \geq 0$.
\end{solution}

\begin{solution}{exo:b2:quadratic:12}
\begin{enumerate}
  \item $\varphi(M, N) = \operatorname{tr}(MN)$ bersifat bilinear dan
        \omterm{def:b2:quadratic:adjoint}{simetrik} (sebab $\operatorname{tr}(MN) =
        \operatorname{tr}(NM)$), dan $\varphi(M, M) = q(M)$:
        jadi $q$ adalah \omterm{def:b2:quadratic:def}{bentuk kuadratik} milik $\varphi$.
  \item Untuk $S$ \omterm{def:b2:quadratic:adjoint}{simetrik} dan $K$ antisimetrik:
        $\operatorname{tr}(SK) =
        \operatorname{tr}\bigl((SK)^{\mathsf T}\bigr) =
        \operatorname{tr}(K^{\mathsf T}S^{\mathsf T}) =
        -\operatorname{tr}(KS) = -\operatorname{tr}(SK)$, jadi
        $\varphi(S, K) = 0$: sehingga kedua subruangnya bersifat
        $\varphi$-ortogonal. Entri demi entri, $\operatorname{tr}
        (M^2) = \sum_{i,j} m_{ij}m_{ji}$: untuk $M$ \omterm{def:b2:quadratic:adjoint}{simetrik}
        ini bernilai $\sum m_{ij}^2 > 0$ (bila $M \neq 0$); sedangkan untuk
        $M$ antisimetrik ia bernilai $-\sum m_{ij}^2 < 0$.
  \item $\mathcal M_n(\R) = S_n \oplus A_n$ dengan dimensi
        $\frac{n(n+1)}2$ dan $\frac{n(n-1)}2$; jadi sebuah \omterm{thm:b2:quadratic:gauss}{reduksi
        Gauss} yang diselaraskan dengan pemecahan
        $\varphi$-ortogonal ini menuliskan $q$ sebagai $\frac{n(n+1)}2$
        kuadrat positif dan $\frac{n(n-1)}2$ kuadrat negatif: jadi signaturnya
        $\bigl(\frac{n(n+1)}2, \frac{n(n-1)}2\bigr)$
        (Sylvester) dan ranknya $n^2$: sehingga bentuknya tak merosot.
\end{enumerate}
\end{solution}

\begin{solution}{pb:b2:quadratic:1}
\textbf{1.} Matriks $G$ \omterm{def:b2:quadratic:adjoint}{simetrik} berkat kesimetrikan hasil kali dalamnya,
dan
\[
X^{\mathsf T}GX = \sum_{i,j}x_ix_j\langle v_i, v_j\rangle
= \Bigl\|\sum_i x_iv_i\Bigr\|^2 \geq 0 ,
\]
dengan kesamaannya bila dan hanya bila $\sum x_iv_i = 0$: jadi $G$ definit
bila dan hanya bila satu-satunya kombinasi nolnya trivial, yakni bila dan hanya
bila keluarganya bebas.

\textbf{2.} Dengan $B = \sqrt A$ (\cref{exo:b2:quadratic:6}): $A
= B^2 = B^{\mathsf T}B$, yakni matriks Gram kolom-kolom $B$;
jadi ambillah $C = B$. Dan $X^{\mathsf T}AX = \norm{CX}^2$, jadi $A$
definit bila dan hanya bila $CX \neq 0$ untuk $X \neq 0$, yakni bila $C$ terbalikkan.

\textbf{3.} Pembatasan bentuknya ke
$\operatorname{Vect}(e_i, e_j)$ bermatriks
$\begin{pmatrix} a_{ii} & a_{ij}\\ a_{ij} & a_{jj}
\end{pmatrix}$, yang masih semidefinit positif: jadi \omterm{def:b2:linalg:det}{determinannya}
(yakni hasil kali \omterm{def:b2:reduction:eigen}{nilai eigennya} yang taknegatif) bernilai $\geq 0$, sehingga
$a_{ij}^2 \leq a_{ii}a_{jj}$. Inilah Cauchy--Schwarz bagi
vektor $v_i, v_j$ pada sebuah perwujudan Gram.

\textbf{4.} \emph{Keberadaannya:} tulislah $A$ sebagai matriks Gram
sebuah keluarga bebas $(v_1, \dots, v_n)$ (pertanyaan 1--2).
Gram--Schmidt menghasilkan $(e_1, \dots, e_n)$ yang ortonormal dengan
\[
v_k = \sum_{i \leq k} t_{ik}\,e_i,
\qquad t_{kk} = \bigl\| v_k - \operatorname{proj}_{k-1}v_k
\bigr\| > 0 ,
\]
jadi $T = (t_{ik})$ segitiga atas dengan diagonal positif,
dan
\[
a_{jk} = \langle v_j, v_k\rangle
= \sum_i t_{ij}t_{ik} = (T^{\mathsf T}T)_{jk} .
\]
\emph{Ketunggalannya:} bila $T_1^{\mathsf T}T_1 = T_2^{\mathsf T}T_2$
maka $R = T_1T_2^{-1}$ memenuhi $R^{\mathsf T}R = I$: jadi $R$
ortogonal, sekaligus segitiga atas dengan diagonal positif
(sebab hasil kali yang demikian). Maka $R^{-1} = R^{\mathsf T}$ bersifat
segitiga atas (sebab balikan yang atas) sekaligus bawah (sebab \omterm{def:b2:linalg:transpose}{transpos} yang
atas): jadi diagonal; dan matriks diagonal yang ortogonal punya
entri $\pm1$, lalu kepositifannya memaksa $R = I$: sehingga $T_1 = T_2$.

\textbf{5.} Untuk $i, j \leq k$, suku $(T^{\mathsf T}T)_{ij} = \sum_m
t_{mi}t_{mj}$ hanya melibatkan $m \leq \min(i,j) \leq k$: jadi
blok terdepan $k\times k$ milik $A$ adalah $T_k^{\mathsf T}T_k$ dengan
$T_k$ blok terdepan $T$. Karena itu $\Delta_k = (\det T_k)^2
= (t_{11}\cdots t_{kk})^2$. Kini \omterm{thm:b2:quadratic:gauss}{reduksi Gauss} sebuah
bentuk yang definit positif, menurut urutan alaminya, tak bertemu koefisien
kuadrat yang nol (sebab pivotnya adalah entri diagonal blok
definit positif yang tereduksi berturut-turut): jadi ia menghasilkan $q
= \sum_k d_k\ell_k^2$ dengan $\ell_k = x_k + (\text{suku dalam }
x_{k+1}, \dots)$, yakni $A = L^{\mathsf T}DL$ dengan $L$
segitiga unipoten; lalu $T = \sqrt D\,L$ merupakan faktor
Cholesky, jadi berkat ketunggalannya $t_{kk}^2 = d_k$ dan
\[
d_k = \frac{(t_{11}\cdots t_{kk})^2}
{(t_{11}\cdots t_{k-1,k-1})^2}
= \frac{\Delta_k}{\Delta_{k-1}} .
\]
Pada \cref{ex:b2:quadratic:tworoads}: $\Delta_1, \Delta_2,
\Delta_3 = 2, 3, 4$ dan pivotnya tadi $2, \frac32, \frac43$.

\textbf{6.} Tiap $a_{ii} = e_i^{\mathsf T}Ae_i > 0$. Misalkan $D =
\operatorname{diag}(a_{ii}^{-1/2})$ dan $B = DAD$: ia definit
positif (lewat \omterm{def:b2:quadratic:def}{kekongruenan}), dengan $b_{ii} = 1$, jadi $\operatorname{tr}
B = n$. \omterm{def:b2:reduction:eigen}{Nilai eigennya} $\mu_i > 0$ memenuhi, menurut AM--GM,
\[
\det B = \prod_i\mu_i
\leq \Bigl(\frac{\sum\mu_i}{n}\Bigr)^{\!n} = 1 ,
\]
dan $\det B = (\det D)^2\det A = \dfrac{\det A}{\prod
a_{ii}}$: jadi $\det A \leq \prod a_{ii}$.

\textbf{7.} AM--GM menjadi kesamaan bila dan hanya bila semua $\mu_i$ sama
(dengan $1$); sedangkan matriks \omterm{def:b2:quadratic:adjoint}{simetrik} yang satu-satunya \omterm{def:b2:reduction:eigen}{nilai eigennya} $1$
adalah $PIP^{\mathsf T} = I$. Jadi kesamaannya berlaku bila dan hanya bila $B = I$,
yakni bila $a_{ij} = 0$ untuk $i \neq j$: sehingga $A$ diagonal.

\textbf{8.} Bila $M$ singular maka kedua ruasnya $\geq 0 =
\abs{\det M}$. Selain itu $A = M^{\mathsf T}M$ bersifat definit
positif dengan $a_{ii} = \norm{c_i}^2$ dan $\det A = (\det
M)^2$: jadi pertanyaan 6 memberi $(\det M)^2 \leq \prod\norm{c_i}^2$.
Kesamaannya berlaku bila dan hanya bila $A = M^{\mathsf T}M$ diagonal (pertanyaan 7),
yakni bila kolomnya ortogonal berpasangan.

\textbf{9.} Besaran $\abs{\det M}$ adalah volume paralelepipedum
yang direntang kolomnya: jadi volumenya paling banyak sebesar hasil kali
panjang rusuknya, dengan kesamaan persis untuk kotak siku-siku.
Bila $\abs{m_{ij}} \leq 1$ maka $\norm{c_i} \leq \sqrt n$, jadi
$\abs{\det M} \leq n^{n/2}$. (Mencapainya memaksa kolomnya
ortogonal berentri $\pm1$: yakni sebuah matriks Hadamard.)

\textbf{10.} Berlaku $X^{\mathsf T}A^{\mathsf T}AX = \norm{AX}^2 > 0$
untuk $X \neq 0$ (sebab $A$ terbalikkan): jadi $A^{\mathsf T}A$ definit
positif. Akar kuadratnya $S$ bersifat definit positif (sebab \omterm{def:b2:reduction:eigen}{nilai eigennya}
$\sqrt{\lambda_i} > 0$), sehingga terbalikkan, lalu $Q = AS^{-1}$
memenuhi
\[
Q^{\mathsf T}Q = S^{-1}A^{\mathsf T}AS^{-1}
= S^{-1}S^2S^{-1} = I :
\]
jadi $A = QS$ dengan $Q$ ortogonal dan $S$ definit positif.

\textbf{11.} Bila $A = QS = Q'S'$ maka $S'^{\,2} =
S'^{\mathsf T}Q'^{\mathsf T}Q'S' = A^{\mathsf T}A = S^2$; sedangkan dua
matriks semidefinit positif yang kuadratnya sama pastilah berimpit
(\cref{exo:b2:quadratic:6}): jadi $S' = S$, lalu $Q' = AS^{-1} =
Q$.

\textbf{12.} Besaran $\det(A + \varepsilon I)$ adalah polinomial tak nol
dalam $\varepsilon$: jadi akarnya berhingga banyak, sehingga suatu barisan
$\varepsilon_k \to 0$ menghindarinya. Tulislah $A + \varepsilon_kI =
Q_kS_k$ (pertanyaan 10). Karena $O(n)$ \omterm{def:b2:metric:compact}{kompak}
(\cref{exo:b2:quadratic:5}), sebuah subbarisan memberi $Q_{\varphi(k)}
\to Q \in O(n)$. Maka
\[
S_{\varphi(k)} = Q_{\varphi(k)}^{\mathsf T}
\bigl(A + \varepsilon_{\varphi(k)}I\bigr)
\longrightarrow Q^{\mathsf T}A =: S,
\]
yang \omterm{def:b2:quadratic:adjoint}{simetrik} semidefinit positif sebagai limit yang demikian (sebab
syaratnya tertutup), dan $A = QS$. Lebih jauh $S^2 = S^{\mathsf T}S =
A^{\mathsf T}QQ^{\mathsf T}A = A^{\mathsf T}A$, jadi $S =
\sqrt{A^{\mathsf T}A}$ berkat ketunggalannya. Untuk $A$ yang singular, $S$
juga singular dan $Q$ tak tunggal: sebab ia dapat diubah sekehendak hati
pada $(\operatorname{im} S)^{\perp}$ --- dengan kasus ekstremnya $A = 0$,
tempat setiap $Q$ yang ortogonal berhasil.

\textbf{13.} Diagonalkan $S = P\Sigma P^{\mathsf T}$ (lewat teorema
spektral), dengan $\Sigma = \operatorname{diag}(\sigma_i)$ dan
$\sigma_i \geq 0$ menyatakan \omterm{def:b2:reduction:eigen}{nilai eigen} $S =
\sqrt{A^{\mathsf T}A}$. Maka
\[
A = QS = (QP)\,\Sigma\,P^{\mathsf T} = U\Sigma V^{\mathsf T},
\qquad U = QP,\ V = P \in O(n) .
\]

\textbf{14.} Berlaku $\norm{Ax}^2 = x^{\mathsf T}S^2x \leq
\sigma_{\max}^2\norm x^2$ dengan kesamaan pada \omterm{def:b2:reduction:eigen}{vektor eigen} teratas
$S$: jadi $\vertiii A_2 = \sigma_{\max}$ --- dan untuk $A$ \omterm{def:b2:quadratic:adjoint}{simetrik},
$S = \sqrt{A^2}$ bernilai eigen $\abs{\lambda_i}$, sehingga memulihkan
\cref{exo:b2:quadratic:7}. \omterm{def:b2:linalg:det}{Determinannya}: $\abs{\det A} =
\abs{\det U}\det\Sigma\abs{\det V} = \sigma_1\cdots\sigma_n$.
Bolanya: dengan menulis $x = Vy$ dan $\norm y = 1$, vektor $Ax = U\Sigma y$
berkoordinat $z_i = \sigma_iy_i$ dalam kerangka ortonormal
kolom-kolom $U$: jadi petanya adalah $\{\sum z_i^2/\sigma_i^2 =
1\}$, yakni sebuah elipsoid dengan setengah sumbu $\sigma_i$.

\textbf{15.} Sebuah perputaran ke sumbu utamanya
(\cref{cor:b2:quadratic:principalaxes}) mengubah $f$ menjadi
$\lambda_1X^2 + \lambda_2Y^2 + \beta_1X + \beta_2Y + c$, dengan
$\lambda_1\lambda_2 = \det A$. Bila $\det A \neq 0$, serap
suku linearnya lewat penggeseran $X \mapsto X -
\frac{\beta_1}{2\lambda_1}$ (dan demikian pula $Y$): jadi $\lambda_1X'^2
+ \lambda_2Y'^2 = c'$. Untuk $\det A > 0$ (yakni tanda yang sama): ia
elips (bila $c'$ bertanda tepat), sebuah titik, atau kosong. Untuk $\det
A < 0$: ia hiperbola (bila $c' \neq 0$) atau dua garis yang bersilangan. Bila
$\det A = 0$ dengan rank $1$ (katakanlah $\lambda_2 = 0 \neq
\lambda_1$): ia $\lambda_1X'^2 + \beta_2Y + c''$, yakni parabola bila
$\beta_2 \neq 0$; dan selain itu dua garis sejajar, satu garis, atau
kosong. Jadi bentuk tak merosotnya: elips, hiperbola, parabola,
yang dikendalikan oleh tanda $\det A$.

\textbf{16.} Bagian kuadratiknya $x^2 + 4xy + y^2$ bermatriks
$\begin{pmatrix}1 & 2\\ 2 & 1\end{pmatrix}$, bernilai eigen $3$
dan $-1$ dengan arah ortonormal $\frac{1}{\sqrt2}(1,1)$ dan
$\frac{1}{\sqrt2}(1,-1)$. Dalam koordinat terputar $u =
\frac{x+y}{\sqrt2}$, $v = \frac{x-y}{\sqrt2}$: berlaku $x^2 + y^2 = u^2
+ v^2$ dan $2xy = u^2 - v^2$, jadi bentuknya $3u^2 - v^2$, dan $2x
- 2y = 2\sqrt2\,v$. Persamaannya menjadi
\[
3u^2 - v^2 + 2\sqrt2\,v = 4
\quad\Longleftrightarrow\quad
3u^2 - \bigl(v - \sqrt2\bigr)^2 = 2 :
\]
yakni sebuah hiperbola, berpusat di $(u, v) = (0, \sqrt2)$, yaitu $(x, y) =
(1, -1)$, dengan sumbu sepanjang kerangka terputarnya.

\textbf{17.} Berlaku $f(x) = (x - x_0)^{\mathsf T}A(x - x_0) + f(x_0)$
setiap kali $Ax_0 = -\frac b2$, yakni $x_0 = -\frac12A^{-1}b$: sebab
gradien $f$ lenyap persis di sana ($x_0$ adalah pusat
kesimetrikannya). Dengan $M = \begin{pmatrix} I & x_0\\ 0 &
1\end{pmatrix}$:
\[
M^{\mathsf T}\widetilde QM
= \begin{pmatrix}
A & Ax_0 + \frac b2\\[2pt]
\bigl(Ax_0 + \frac b2\bigr)^{\mathsf T} &
x_0^{\mathsf T}Ax_0 + b^{\mathsf T}x_0 + c
\end{pmatrix}
= \begin{pmatrix} A & 0\\ 0 & f(x_0)\end{pmatrix},
\]
dan $\det M = 1$: jadi $\det\widetilde Q = f(x_0)\det A$. Persamaan
terpusatnya berbunyi $q(X) = -f(x_0)$: untuk $f(x_0) = 0$ ia
merosot menjadi $q(X) = 0$ (yakni dua garis lewat pusatnya bila
signaturnya $(1,1)$, atau titik tunggal $x_0$ bila $q$
definit); sedangkan untuk $f(x_0) \neq 0$ koniknya berupa elips
atau hiperbola yang sejati.

\textbf{18.} Berlaku $A^{-1} = -\frac13\begin{pmatrix} 1 & -2\\ -2 &
1\end{pmatrix}$ dan $\frac b2 = (1, -1)$, jadi $x_0 =
-A^{-1}\frac b2 = (1, -1)$, seperti ditemukan pada pertanyaan 16. Lalu $f(x_0) =
q(1,-1) + 2 + 2 - 4 = (1 - 4 + 1) + 0 = -2 \neq 0$, dan
$\det\widetilde Q = f(x_0)\det A = (-2)(-3) = 6 \neq 0$: jadi
tak merosot; dan $\det A = -3 < 0$: sehingga sebuah hiperbola --- dan memang
persamaan terpusatnya $3u^2 - (v - \sqrt2)^2 = -f(x_0) = 2$
cocok dengan pertanyaan 16.

\textbf{19.} Matriksnya adalah $(1-\lambda)I + \lambda J$, yang
bernilai eigen $1 + 2\lambda$ (pada arah $(1,1,1)$) dan $1 -
\lambda$ (ganda, pada $x + y + z = 0$). Kasusnya:
\begin{itemize}
  \item $-\frac12 < \lambda < 1$: semua \omterm{def:b2:reduction:eigen}{nilai eigennya} positif, jadi
        sebuah elipsoid putaran terhadap $(1,1,1)$ (dan sebuah bola
        bila $\lambda = 0$);
  \item $\lambda = 1$: $q = (x+y+z)^2$, jadi persamaannya memberi
        dua bidang sejajar $x + y + z = \pm1$;
  \item $\lambda = -\frac12$: \omterm{def:b2:reduction:eigen}{nilai eigennya} $0, \frac32,
        \frac32$, jadi sebuah silinder lingkaran bersumbu $(1,1,1)$;
  \item $\lambda > 1$: signaturnya $(1, 2)$, jadi sebuah hiperboloid
        berdaun dua;
  \item $\lambda < -\frac12$: signaturnya $(2, 1)$, jadi sebuah
        hiperboloid berdaun satu.
\end{itemize}

\textbf{20.} Bentuk polar $q$ adalah sebuah hasil kali dalam
$\langle\cdot,\cdot\rangle_q$ pada $E$. Untuk $x$ yang tetap, pemetaan $y \mapsto
\varphi'(x, y)$ (yakni bentuk polar $q'$) bersifat linear, sehingga sama dengan
$\langle z_x, y\rangle_q$ bagi $z_x$ yang tunggal; lalu $u(x) := z_x$ bersifat
linear (berkat ketunggalannya), dan $\langle u(x), y\rangle_q =
\varphi'(x,y) = \varphi'(y,x) = \langle u(y), x\rangle_q$: jadi $u$
\omterm{def:b2:quadratic:adjoint}{simetrik} dalam ruang Euklides $(E,
\langle\cdot,\cdot\rangle_q)$. Teorema spektral
(\cref{thm:b2:quadratic:spectral}) memberi basis eigen yang
$q$-ortonormal, sebut $(\varepsilon_i)$, dengan $u(\varepsilon_i) =
\mu_i\varepsilon_i$: jadi di dalamnya $q(x) = \sum x_i^2$ dan $q'(x) =
\langle u(x), x\rangle_q = \sum\mu_ix_i^2$.

\textbf{21.} Misalkan $P$ matriks basis itu: maka \omterm{def:b2:quadratic:def}{kekongruenannya}
memberi $P^{\mathsf T}AP = I$ dan $P^{\mathsf T}BP =
\operatorname{diag}(\mu_i)$. Lalu
\[
\det(B - \mu A) = \det(P^{-\mathsf T})
\det\bigl(\operatorname{diag}(\mu_i) - \mu I\bigr)
\det(P^{-1})
= (\det P)^{-2}\prod_i(\mu_i - \mu) :
\]
jadi $\mu_i$ adalah akar berkas $\det(B - \mu A)$.

\textbf{22.} Berlaku $\det(B - \mu A) =
\det\begin{pmatrix} -2\mu & 1-\mu\\ 1-\mu & -\mu\end{pmatrix}
= 2\mu^2 - (1-\mu)^2 = \mu^2 + 2\mu - 1$, jadi akarnya $\mu_\pm = -1
\pm \sqrt2$. Menyelesaikan $(B - \mu_\pm A)v = 0$ memberi $v_\pm = (1 -
\mu_\pm,\ 2\mu_\pm)$ (dan kesamaan baris keduanya $(1-\mu)^2 =
2\mu^2$ pada akarnya menegaskannya). Norma-$A$-nya keluar
rapi: $q_A(v_\pm) = 2(1 + \mu_\pm^2)$, lalu kita periksa
$\varphi_A(v_+, v_-) = 0$ dengan memakai $\mu_+ + \mu_- = -2$ dan
$\mu_+\mu_- = -1$. Jadi basis $\Bigl(\frac{v_+}{\sqrt{2(1 +
\mu_+^2)}}, \frac{v_-}{\sqrt{2(1+\mu_-^2)}}\Bigr)$ bersifat
ortonormal terhadap $A$ dan mendiagonalkan $B$ dengan entri
$\mu_\pm$.

\textbf{23.} Andaikan suatu $P$ terbalikkan mendiagonalkan kedua bentuknya,
maka perhitungan pertanyaan 21 akan memberi $\det(B - \mu A) = (\det
P)^{-2}\prod(d_{2i} - \mu d_{1i})$, yakni polinomial real yang terurai
menjadi faktor linear yang real. Tetapi di sini
\[
\det(B - \mu A) = \det\begin{pmatrix} -\mu & 1\\ 1 &
\mu\end{pmatrix} = -\mu^2 - 1 ,
\]
yang berderajat $2$ tanpa akar real: jadi kontradiksi. (Sebab
$A$ bersignatur Lorentz mengizinkan ``perputaran''-$B$ tanpa sumbu
yang real.)

\textbf{24.} Berlaku $A^{-1}B = A^{-1/2}\bigl(A^{-1/2}BA^{-1/2}\bigr)
A^{1/2}$ dengan $A^{1/2} = \sqrt A$ yang definit positif
(\cref{exo:b2:quadratic:6}): jadi $A^{-1}B$ serupa dengan
matriks \emph{\omterm{def:b2:quadratic:adjoint}{simetrik}} $A^{-1/2}BA^{-1/2}$, sehingga \omterm{def:b2:reduction:diag}{dapat didiagonalkan} dengan
\omterm{def:b2:reduction:eigen}{nilai eigen} yang real. Dan $\det(B - \mu A) = \det A\cdot
\det(A^{-1}B - \mu I)$: jadi akar berkas $\mu_i$ pada pertanyaan
21 persis merupakan \omterm{def:b2:reduction:eigen}{nilai eigen} $A^{-1}B$.

\textbf{25.} (i) Setiap Bagiannya menyandari teorema spektral:
lewat akar kuadratnya (Cholesky, polar), batas \omterm{def:b2:reduction:eigen}{nilai
eigennya} (Hadamard), sumbu utamanya (konik), dan
versi yang $q$-selaras (reduksi serentak). (ii) Perwujudan
Gram, Hadamard, dan \omterm{pb:b2:quadratic:1}{dekomposisi polar} bertahan di
dunia yang semidefinit; sedangkan ketunggalan Cholesky, rumus
pivotnya, dan reduksi serentaknya menuntut kedefinitan
(dan pertanyaan 12 serta 23 menunjukkan persis bagaimana keduanya gagal). (iii)
Osilasi kecil yang terkopel: energi kinetiknya (yang definit positif)
dan energi potensialnya adalah dua \omterm{def:b2:quadratic:def}{bentuk kuadratik}; lalu basis
pertanyaan 20--22 adalah \emph{moda normal} sistemnya, dan
$\mu_i$ adalah kuadrat frekuensi sudutnya.

\end{solution}
